\documentclass[11pt]{article}
\usepackage[paperwidth=7.10in,paperheight=7.75in,left=.43in,right=.43in,top=.28in,bottom=.42in,footskip=.19in]{geometry}
\usepackage{amsmath}
\usepackage[hidelinks]{hyperref}
\pdfinfoomitdate=1
\pdftrailerid{}
\hypersetup{pdfauthor={Kornstein Lab for Antisemitism Research},pdfsubject={Research proposal; illustrative methods, no empirical findings}}
\setlength{\parindent}{0pt}
\setlength{\parskip}{3pt}
\setlength{\abovedisplayskip}{7pt}
\setlength{\belowdisplayskip}{7pt}
\pagestyle{plain}
\makeatletter
\renewcommand\section{\@startsection{section}{1}{0pt}{9pt plus 1pt minus 1pt}{3pt}{\normalfont\large\bfseries}}
\makeatother
\newcommand{\papertitle}[1]{%
  \hypersetup{pdftitle={#1}}%
  \begin{center}
    {\fontsize{18.5}{22}\selectfont #1\par}
    \vspace{6pt}
    {\fontsize{13}{16}\selectfont Kornstein Lab for Antisemitism Research\par}
    {\fontsize{12}{15}\selectfont Research proposal\par}
  \end{center}
  \vspace{2pt}
}
\newenvironment{paperabstract}{\begin{list}{}{\leftmargin=21pt\rightmargin=21pt}\item[]\begin{center}\textbf{Abstract}\end{center}\vspace{-3pt}\small}{\end{list}\vspace{2pt}}
\newenvironment{methodlist}{\begin{enumerate}\setlength{\itemsep}{1pt}\setlength{\parsep}{0pt}\setlength{\topsep}{2pt}}{\end{enumerate}}

\begin{document}
\papertitle{Mapping the spread of antisemitic narratives}
\begin{paperabstract}
This proposed study examines how antisemitic narratives recur and change across online discourse. It combines text analysis and network methods while distinguishing observed connections from claims about influence.
\end{paperabstract}

\section{Introduction}
We propose to examine recurring narratives and their temporal relationships across sampled online discourse. The aim is to describe observable patterns without treating similarity as evidence of coordination.

\section{Network framework}
An illustrative thresholded similarity network is
\[
  A_{ij}(t)=\mathbf{1}\!\left\{s_{ij}(t)\geq\tau\right\},
  \qquad i\neq j.
\]
Here, $s_{ij}(t)$ denotes a specified similarity measure at time $t$ and $\tau$ a chosen threshold. An edge represents similarity under this definition; it does not establish influence or causal transmission.

\section{Proposed methodology}
\begin{methodlist}
  \item Corpus construction.
  \item Narrative identification.
  \item Temporal network analysis.
\end{methodlist}

\section{Limitations}
Missing observations, repeated quotations, and threshold choices may alter network structure. Sensitivity analysis and contextual review would be required. This proposal reports no empirical network findings.
\end{document}
